\chapter{Literatür Özeti}
\label{ch:literatur}

Bu bölümde, mevcut çalışmanın temellendiği dört eksende literatür
değerlendirilmiştir: (i) yüz tanıma yöntemleri ve performans metrikleri,
(ii) gömülü platformlarda (MCU/SBC) gerçek zamanlı uygulama kısıtları,
(iii) IoT tabanlı acil bildirim sistemleri ve (iv) biyometrik verilerin
yasal/etik çerçevesi.

\section{Yüz Tanıma Yöntemleri ve Performans Metrikleri}

Modern yüz tanıma; tespit (\emph{detection}), hizalama (\emph{alignment})
ve özellik çıkarımı (\emph{embedding}) olmak üzere üç ardışık adımdan
oluşmaktadır. ArcFace mimarisi, son katmandaki softmax fonksiyonuna eklenen
\emph{additive angular margin} terimi sayesinde, sınıflar arası ayrımı
kosinüs uzayında belirginleştirir ve LFW, MegaFace gibi standart veri
kümelerinde insan performansına yakın doğruluk üretir \cite{deng2019arcface}.
ArcFace'in açık kaynak referans uygulaması InsightFace
\cite{guo2021insightface}, RetinaFace \cite{deng2020retinaface} ile
birleştirilmiş hâlde \texttt{buffalo\_l} model paketi olarak
yayınlanmaktadır; bu paket bu çalışmanın bulut katmanında kullanılmıştır.

Yüz tanıma sistemlerinin performansı tipik olarak iki tamamlayıcı oranla
nicelendirilir:

\begin{itemize}
    \item \textbf{Yanlış Kabul Oranı (FAR --- False Accept Rate):} aranan
          listede bulunmayan bir kişinin, aranan biri ile eşleştiği vakaların
          oranı.
    \item \textbf{Yanlış Reddetme Oranı (FRR --- False Reject Rate):}
          gerçek bir aranan kişinin sistem tarafından kaçırıldığı vakaların
          oranı.
\end{itemize}

Bu iki oran, sistemin karar eşiği (kosinüs benzerliği için $\tau$) bir
parametre olarak süpürüldüğünde \emph{karşılıklı ödünleşim} (ROC eğrisi)
şeklinde gözlemlenir. İki oranın eşitlendiği nokta \textbf{Eşit Hata
Oranı (EER --- Equal Error Rate)} olarak adlandırılır ve sistemin tek
sayısal kalite göstergesidir.

Tek bir karelik snapshot ile gerçekleştirilen tanıma, hareket bulanıklığı,
ani aydınlatma değişimi ve geçici tıkanmalardan etkilenir. Birden çok
kare üzerinde gömme vektörlerinin centroid'ini hesaplama yaklaşımı,
literatürde yaygın olarak kullanılan ve tek-shot yaklaşıma kıyasla
gözlemlenebilir doğruluk artışı sağlayan bir tekniktir
\cite{chen2018multiframe}. Bu çalışmada da 10 karelik öbeklenme tercih
edilmiştir.

\section{Sahtelik ve Canlılık (Liveness)}

Yüz tanıma sistemlerinin pratik bir güvenlik açığı, basılı fotoğraf veya
ekran tekrarı (\emph{photo replay}) saldırılarıdır. \emph{Aktif canlılık}
yöntemleri (kullanıcıdan gülümsemesi, gözünü kırpması istenmesi)
caydırıcılığı yüksektir ancak kullanıcı etkileşimi gerektirir; bu da
sürücü için bir dikkat dağıtıcıdır. Bu nedenle, çoklu kare gömme
vektörlerinin kareler arası standart sapması (\emph{cross-frame standard
deviation}) bir \textbf{pasif canlılık} göstergesi olarak kullanılmıştır
\cite{tang2018liveness}. Gerçek bir canlı yüz, mikro-hareketler ve doğal
yüz kasları nedeniyle düşük ama sıfırdan farklı bir standart sapma
üretirken; bir fotoğraf replay senaryosunda standart sapma anlamlı
biçimde azalır. Bu sinyal ek bir model yükü getirmeden, mevcut gömme
vektörleri üzerinden ücretsiz olarak hesaplanmaktadır.

\section{Gömülü Platformlarda Gerçek Zamanlı Uygulama Kısıtları}

Yüz tanıma iş yükünün gömülü tarafa yerleştirilmesi, üç temel kısıt
nedeniyle bu projede tercih edilmemiştir:

\begin{enumerate}
    \item \textbf{Hesaplama yükü:} ArcFace-R100 yaklaşık 24 milyon
          parametreye sahip bir CNN'dir. INT8 nicemleme sonrası bile
          STM32F767 sınıfı bir MCU için (240 MHz, $\leq$~512~KB RAM)
          gecikme bütçesi aşılmaktadır.
    \item \textbf{Bellek kısıtı:} buffalo\_l model paketi, INT8'e
          indirilse dahi 30~MB'ı aşmaktadır; Cortex-M7 MCU'larında
          bu kalıcı bellek sığasının ötesindedir.
    \item \textbf{Geliştirme döngüsü:} Sergi süresi (12 gün) gibi sınırlı
          bir takvimde, hassas nicemleme ve donanım hızlandırma
          çalışmaları öncelikten düşmektedir.
\end{enumerate}

Bu kısıtlar, MCU üzerinde yalnızca \emph{olay yönetimi} (state machine,
LED/buzzer/buton, BLE köprüsü) çalıştırma; tanımayı sunucuya devretme
kararını desteklemektedir. Aynı yaklaşım, Raspberry Pi 4 gibi
tek-kart-bilgisayar (SBC) tabanlı çalışmalarda da
``edge inference + cloud verification'' hibridi olarak
gözlemlenmektedir \cite{kumar2021edgeface}.

\section{IoT Tabanlı Acil Bildirim Sistemleri}

Acil durum tetikleme bileşeni için iki yaygın yaklaşım vardır:
(i) doğrudan gömülü cihazdan GSM modülü ile arama/SMS, (ii) BLE üzerinden
eşli akıllı telefon aracılığı ile arama. İlk yaklaşım, donanım maliyeti
(SIM800/SIM7600 modülleri), SIM kart yönetimi ve operatör onayı gibi ek
yükler getirir. İkinci yaklaşım ise sürücünün zaten yanında bulunan akıllı
telefonu \emph{ağ geçidi} (gateway) olarak kullanır; sistem karmaşıklığını
ve maliyetini düşürür. Bu çalışma ikinci yaklaşımı benimsemiştir.

İki büyük mobil platformun bu konudaki tutumları farklıdır: Android,
\texttt{Intent.ACTION\_CALL("tel:155")} ile arka planda otomatik çağrı
başlatılmasına izin verir \cite{android2024telephony}; iOS ise üçüncü
taraf uygulamalardan tam-otomatik aramaya izin vermez, yalnızca
\texttt{tel://}<numara> URL şemasıyla numara yüklenmiş dialer ekranını
açar \cite{apple_tel_url}. Mevcut tez için, ekibin elinde bulunan iPhone
ile sergiye gidilebilmesi pratik bir tercih olarak iOS'un seçilmesini
sağlamıştır; kullanıcının dialer'da yeşil tuşa tek dokunuşu, hem Apple
sandbox gereği hem de \emph{yanlış pozitif arama koruması} olarak
savunulan bir adım hâline gelmiştir. Android tarafında aynı sistem,
\texttt{Intent.ACTION\_CALL} ile tek-dokunuşsuz çalışacak şekilde
genişletilebilir; bu olasılık \ref{ch:sonuc-oneriler}'te ele alınmıştır.

\section{Biyometrik Verilerin Yasal ve Etik Çerçevesi}

Yüz görüntüsü, KVKK kapsamında \emph{özel nitelikli kişisel veri} olarak
sınıflandırılmaktadır \cite{kvkk2016}. Bu kapsamda toplanan verinin:
(i) açık rıza ile veya kanunlarda öngörülmüş hâllerde işlenmesi,
(ii) işlemenin amacına uygun, asgari ölçüde toplanması,
(iii) öngörülen süre boyunca tutularak sonrasında silinmesi
gerekmektedir. Mevcut tasarımda yolcu görüntüleri yalnızca işleme süresince
RAM'de tutulmakta, sunucunun diskine yazılmamaktadır; sonuç dönüşünden
sonra ilgili oturum yapısı bellekten silinmektedir. Aranan kişiler veri
tabanına yapılacak girişlerin ise yetkili idari/kolluk birimleri
tarafından sağlanması, sistem mimarisinin bir varsayımıdır.

GDPR'nin \emph{purpose limitation} ilkesi de \cite{gdpr2018}, sistemin
yalnızca aranan kişi tespiti amacıyla çalışabileceğini; aynı altyapının
müşteri profillemesi veya benzeri amaçlarla kullanılamayacağını
gerektirir. Tasarımda saklanan tek kalıcı veri, aranan kişilere ait
embedding vektörleridir; ham fotoğraf veri tabanına yazılmaz.

\section{Bu Çalışmanın Konumu}

Yukarıdaki dört eksen değerlendirildiğinde, mevcut çalışmanın özgün
katkıları şu şekilde özetlenebilir:

\begin{itemize}
    \item Açık kaynak InsightFace tabanlı bir tanıma sunucusunun, ESP32-CAM
          ile birlikte uçtan uca taksi senaryosuna uyarlanması;
    \item Çoklu kare centroid agregasyonu ve cross-frame standart sapması
          ile ek model yükü getirmeyen \emph{pasif canlılık} sinyalinin
          birlikte uygulanması;
    \item Sürücü etkileşimini tek butona indiren ve dikkat dağıtmayan bir
          STM32 olay yöneticisi tasarımı;
    \item Akıllı telefonun BLE ağ geçidi olarak kullanılmasıyla, GSM
          modülü gerektirmeyen düşük maliyetli bir acil çağrı zinciri;
    \item Toplam donanım eki maliyetinin 500~TL altında tutulması ve KVKK
          ilkelerine uygun veri yönetimi.
\end{itemize}
